\begin{afterword}
\summaryhead

The post asks why rejecting a belief because of its causes should count as a fallacy, when the causes of a belief are what make it reliable. Articles on the genetic fallacy, it notes, say that origins sometimes matter, as with a trusted expert, and sometimes do not, as with Kekulé's dream of the benzene ring.

The author's answer is that the fallacy is real in principle, because a belief's original cause is not its current justification. But people change their minds less often than they think, so origins matter more for humans than for ideal reasoners. When a source you relied on is discredited, as the Bible is for someone who stops believing, you should clear your mind and doubt everything that came from it, without simply reversing it. Where clear-cut evidence exists, as for benzene, the source no longer matters.

The genetic fallacy, the post concludes, is a fallacy when other justifications exist but the accusation is presented as settling the issue. It ends with four rules of thumb.

\responsehead

The post reaches the standard account of the genetic fallacy by a long route, and adds one idea of its own.

The route begins by calling the fallacy ``a very strange idea'' and mocking the articles that say origins sometimes matter and sometimes do not. The post's own conclusion is the same. The source stops mattering once there is clear-cut evidence, as with Kekulé's benzene ring, and needs attention when there is none, as with experts. The definition the post ends on, an accusation ``presented as if it settled the issue,'' is what standard accounts express with the word ``solely.''

The idea it adds is sound. People revise beliefs poorly, so a belief's origin tells you more about a human's belief than about an ideal reasoner's, and when a source is discredited you should doubt what you learned from it, without simply reversing it. The support offered is the job-choice survey of ``We Change Our Minds Less Often Than We Think,'' which does not concern discredited sources. The belief-perseverance studies do. They support the claim: impressions persisted after the evidence behind them was discredited. They also complicate the post's remedy. In the first of these studies, explaining the perseverance process to subjects generally removed the mistaken self-perceptions, which sits uneasily with the post's ``convulsive effort.''

The main example is asserted, not shown. People who stopped believing that God wrote the Bible but still value its ethics ``have failed to clear their minds.'' No case is given. Whether any such person re-examined those ethics is a question of fact about that person, and the post's own rules speak of doubting a source's teachings, not of verdicts on the people who kept some of them.

One sentence near the end, crediting Hal Finney with the name ``genetic heuristic,'' comes from a comment on this post and was added after publication.

In short: a long route to the standard account of the genetic fallacy, one sound addition about discredited sources, supported by a survey about something else when better evidence existed, and a main example asserted about unnamed people.
\end{afterword}
